\section{Abdominal X-Ray}


\section*{Patient focus}
This chapter explains the examination or clinical problem in accessible language. Patients should use it to prepare questions and understand the purpose, experience, benefits, and limitations of imaging; local instructions and the official report take priority.


\subsubsection*{What Is This Test?}
An abdominal X-ray images your abdominal organs and structures. It's primarily used to evaluate bowel obstruction, kidney stones, foreign bodies, and perforation (free air).

\subsubsection*{Alternative Names}
AXR, KUB (kidneys, ureters, bladder), plain abdominal radiograph, abdominal radiograph


\subsubsection*{Principle}
Abdominal X-ray uses X-ray photons generated by a tungsten cathode. Photons pass through the body and are attenuated by tissues of different density. A detector (film, CR plate, or digital flat panel) captures the transmitted photons. Air appears black, fat appears dark gray, soft tissue appears light gray, bone appears white, and metal appears bright white.

\subsubsection*{History}
The abdominal X-ray was one of the first clinical applications of X-rays following Wilhelm Röntgen's discovery in 1895. It became a cornerstone of emergency radiology in the early 20th century.

\subsubsection*{Why This Test Is Ordered}
Your doctor orders an abdominal X-ray when you have severe abdominal pain, nausea and vomiting, suspected kidney stones, or when you have had recent surgery and your bowel function needs evaluation. It answers the question: is there a bowel obstruction, perforation, kidney stone, or abnormal gas pattern in your abdomen?

\subsubsection*{Is This a Primary or Secondary Test}
A primary test for initial evaluation of acute abdominal symptoms. More advanced imaging like CT is often performed afterward if the X-ray is inconclusive.

\subsubsection*{Limits of abdominal radiograph accuracy}
A plain abdominal film shows selected patterns, but a normal result does not exclude many causes of obstruction, perforation, or acute abdominal pain. Performance depends on the suspected diagnosis, patient position, and views obtained.

\subsubsection*{How This Test Is Performed}
You lie on a table or stand in front of the X-ray machine. Images are taken of your abdomen from the front. You may be asked to hold your breath briefly. The test takes about fifteen minutes.

\subsubsection*{What Preparation Is Needed}
Usually no preparation is needed. Tell your doctor if you are pregnant.

\subsubsection*{Contraindications and Precautions}
Pregnancy is not an automatic contraindication; the examination should be justified and optimized, and alternatives considered when appropriate. For contrast-enhanced studies, review the indication, prior severe reactions, and renal function according to local protocol.

\subsubsection*{Risks}
Risks associated with X-ray imaging include:
\begin{itemize}
    \item \textbf{Radiation exposure}: Exposure to low levels of ionizing radiation, with cumulative lifetime dose optimization guided by the ALARA (As Low As Reasonably Achievable) principle.
    \item \textbf{Contrast reaction}: Allergic reactions can occur, including rare severe reactions. Kidney injury after iodinated contrast is uncommon and risk depends on baseline kidney function, illness, and the clinical context.
\end{itemize}


Results are typically available within thirty minutes. You can resume normal activities immediately.

\subsubsection*{Understanding Your Results}

\subsubsection*{Normal Results}
A normal result shows normal bowel gas pattern with no obstruction, perforation, or abnormal calcifications.

\subsubsection*{Follow-up Tests}
If the X-ray is abnormal or inconclusive, your doctor may order a CT scan of the abdomen for more detailed evaluation, a contrast-enhanced study to better visualize structures, or an ultrasound to assess specific organs.
\clearpage
